%!TEX root = ../main.tex
% Chapitre 8 : Modeles Generatifs — GAN
\chapter{Mod\`eles G\'en\'eratifs --- GAN}
\label{ch:gan}

\paragraph{Objectifs du chapitre}
\begin{itemize}
\item Distinguer mod\`eles g\'en\'eratifs et discriminatifs
\item D\'eriver l'objectif minimax des GAN et le discriminateur optimal
\item Comprendre les probl\`emes d'entra\^inement et les solutions (WGAN, WGAN-GP)
\item Impl\'ementer un DCGAN complet en PyTorch
\item Conna\^itre les m\'etriques d'\'evaluation (FID, IS)
\end{itemize}

\section{Mod\`eles g\'en\'eratifs vs discriminatifs}
\label{sec:gen-vs-disc}

\begin{definition}[Mod\`ele discriminatif\index{mod\`ele discriminatif}]
Un mod\`ele discriminatif apprend directement la distribution conditionnelle $p(y \mid x)$, c'est-\`a-dire la fronti\`ere de d\'ecision entre les classes. Exemples : r\'egression logistique, SVM, r\'eseaux de neurones classifieurs.
\end{definition}

\begin{definition}[Mod\`ele g\'en\'eratif\index{mod\`ele g\'en\'eratif}]
Un mod\`ele g\'en\'eratif apprend la distribution jointe $p(x, y)$ ou la distribution marginale $p(x)$ des donn\'ees. Il peut \emph{g\'en\'erer} de nouveaux \'echantillons. Exemples : GAN, VAE (cf.\ chapitre~\ref{ch:vae}), mod\`eles de diffusion, mod\`eles autoregressifs.
\end{definition}

\begin{remarque}
Un mod\`ele g\'en\'eratif permet non seulement de cr\'eer de nouvelles donn\'ees, mais aussi d'effectuer du raisonnement probabiliste (d\'etection d'anomalies, compl\'etion, etc.).
\end{remarque}

\section{Cadre GAN : l'objectif minimax}
\label{sec:gan-framework}

\subsection{Formulation}

Un \textbf{Generative Adversarial Network}\index{GAN} (Goodfellow et al., 2014) se compose de deux r\'eseaux en comp\'etition :
\begin{itemize}
    \item Le \textbf{g\'en\'erateur} $G : \mathcal{Z} \to \mathcal{X}$ transforme un bruit al\'eatoire $z \sim p_z(z)$ en un \'echantillon $G(z)$.
    \item Le \textbf{discriminateur} $D : \mathcal{X} \to [0, 1]$ estime la probabilit\'e qu'un \'echantillon provienne des vraies donn\'ees plut\^ot que du g\'en\'erateur.
\end{itemize}

\begin{mathbox}[title=\textbf{Objectif minimax des GAN}]\index{objectif minimax}
\begin{equation}
\min_G \max_D \; V(D, G) = \E_{x \sim p_{\text{data}}}\!\big[\log D(x)\big] + \E_{z \sim p_z}\!\big[\log\!\big(1 - D(G(z))\big)\big]
\label{eq:gan-minimax}
\end{equation}
\end{mathbox}

\subsection{D\'erivation du discriminateur optimal $D^*$}

\begin{theoreme}[Discriminateur optimal\index{discriminateur optimal}]
\label{thm:optimal-D}
Pour un g\'en\'erateur $G$ fix\'e, le discriminateur optimal est :
\begin{equation}
D^*_G(x) = \frac{p_{\text{data}}(x)}{p_{\text{data}}(x) + p_g(x)}
\label{eq:optimal-D}
\end{equation}
o\`u $p_g$ est la distribution induite par $G$.
\end{theoreme}

\begin{proof}
Pour $G$ fix\'e, on maximise $V(D, G)$ par rapport \`a $D$. L'int\'egrande vaut, pour chaque $x$ :
\begin{equation}
f(D(x)) = p_{\text{data}}(x) \log D(x) + p_g(x) \log(1 - D(x))
\end{equation}
C'est une fonction de $D(x) \in [0, 1]$. En d\'erivant et en annulant :
\begin{align}
\frac{\partial f}{\partial D(x)} &= \frac{p_{\text{data}}(x)}{D(x)} - \frac{p_g(x)}{1 - D(x)} = 0 \\
\Rightarrow \quad p_{\text{data}}(x)\big(1 - D(x)\big) &= p_g(x) D(x) \\
\Rightarrow \quad D^*(x) &= \frac{p_{\text{data}}(x)}{p_{\text{data}}(x) + p_g(x)}
\end{align}
On v\'erifie que la d\'eriv\'ee seconde est n\'egative, confirmant le maximum.
\end{proof}

\subsection{D\'erivation de l'optimalit\'e du g\'en\'erateur}

\begin{theoreme}[Optimalit\'e du g\'en\'erateur et divergence de Jensen-Shannon\index{divergence de Jensen-Shannon}]
\label{thm:optimal-G}
En rempla\c{c}ant $D^*$ dans l'objectif, on obtient :
\begin{equation}
C(G) = V(D^*_G, G) = -\log 4 + 2 \cdot \KL{JSD}(p_{\text{data}} \| p_g)
\end{equation}
o\`u la divergence de Jensen-Shannon est :
\begin{equation}
\text{JSD}(p \| q) = \frac{1}{2}\KL(p \| m) + \frac{1}{2}\KL(q \| m), \quad m = \frac{p + q}{2}
\end{equation}
Le minimum global $C(G) = -\log 4$ est atteint si et seulement si $p_g = p_{\text{data}}$.
\end{theoreme}

\begin{proof}
En substituant $D^*$ :
\begin{align}
C(G) &= \E_{x \sim p_{\text{data}}}\!\left[\log \frac{p_{\text{data}}(x)}{p_{\text{data}}(x) + p_g(x)}\right] + \E_{x \sim p_g}\!\left[\log \frac{p_g(x)}{p_{\text{data}}(x) + p_g(x)}\right] \\
&= \E_{p_{\text{data}}}\!\left[\log \frac{p_{\text{data}}}{2 \cdot \frac{p_{\text{data}} + p_g}{2}}\right] + \E_{p_g}\!\left[\log \frac{p_g}{2 \cdot \frac{p_{\text{data}} + p_g}{2}}\right] \\
&= -\log 4 + \KL\!\left(p_{\text{data}} \,\Big\|\, \frac{p_{\text{data}} + p_g}{2}\right) + \KL\!\left(p_g \,\Big\|\, \frac{p_{\text{data}} + p_g}{2}\right) \\
&= -\log 4 + 2 \cdot \text{JSD}(p_{\text{data}} \| p_g)
\end{align}
Puisque $\text{JSD} \geq 0$ avec \'egalit\'e ssi $p_{\text{data}} = p_g$, le minimum est bien $-\log 4$.
\end{proof}

\subsection{Diagramme TikZ de l'entra\^inement GAN}

\begin{figure}[htbp]
\centering
\begin{tikzpicture}[
    netblock/.style={rectangle, draw, fill=blue!12, minimum width=2.8cm, minimum height=1.2cm, align=center, rounded corners=3pt, font=\small\bfseries},
    datablock/.style={rectangle, draw, fill=green!12, minimum width=2.2cm, minimum height=0.9cm, align=center, rounded corners=3pt, font=\small},
    lossblock/.style={rectangle, draw, fill=red!12, minimum width=2.5cm, minimum height=0.9cm, align=center, rounded corners=3pt, font=\small},
    arr/.style={-{Stealth}, thick},
    darr/.style={-{Stealth}, thick, dashed},
]

% Generateur
\node[datablock] (noise) at (0, 0) {$z \sim \mathcal{N}(0, I)$};
\node[netblock] (G) at (3, 0) {G\'en\'erateur\\$G_\theta$};
\node[datablock] (fake) at (6.5, 0) {$G(z)$\\(faux)};

% Vraies donnees
\node[datablock] (real) at (6.5, 2.5) {$x \sim p_{\text{data}}$\\(vrai)};

% Discriminateur
\node[netblock] (D) at (10, 1.25) {Discriminateur\\$D_\phi$};

% Sortie
\node[lossblock] (loss) at (13.5, 1.25) {$V(D, G)$\\Loss};

% Fleches
\draw[arr] (noise) -- (G);
\draw[arr] (G) -- (fake);
\draw[arr] (fake) -- (D);
\draw[arr] (real) -- (D);
\draw[arr] (D) -- (loss);

% Retropropagation
\draw[darr, red!60!black] (loss.south) -- ++(0, -1.2) -| (D.south)
    node[pos=0.25, below, font=\footnotesize, text=red!60!black] {$\nabla_\phi$ (maximiser $V$)};
\draw[darr, blue!60!black] (loss.north) -- ++(0, 1.0) -| (G.north)
    node[pos=0.25, above, font=\footnotesize, text=blue!60!black] {$\nabla_\theta$ (minimiser $V$)};

\end{tikzpicture}
\caption{Boucle d'entra\^inement d'un GAN. Le discriminateur $D$ est mis \`a jour pour maximiser $V$, tandis que le g\'en\'erateur $G$ est mis \`a jour pour le minimiser. Les fl\`eches en pointill\'e repr\'esentent la r\'etropropagation.}
\label{fig:gan-training}
\end{figure}

\section{Probl\`emes d'entra\^inement}
\label{sec:gan-training-problems}

\subsection{Effondrement de mode (Mode Collapse)}

\begin{definition}[Mode collapse\index{mode collapse}]
Le g\'en\'erateur se <<~sp\'ecialise~>> dans la production d'un petit nombre de modes de la distribution cible, ignorant la diversit\'e des donn\'ees r\'eelles. Formellement, $p_g$ est concentr\'ee sur un sous-ensemble de mesure faible du support de $p_{\text{data}}$.
\end{definition}

\subsection{Gradients \'evanescents pour $G$}

\begin{attention}
Si le discriminateur est trop performant ($D(x) \approx 1$ pour les vraies donn\'ees et $D(G(z)) \approx 0$ pour les fausses), alors $\log(1 - D(G(z))) \approx \log(1) = 0$ et le gradient pour $G$ devient quasi nul.

\textbf{Solution pratique :} au lieu de minimiser $\log(1 - D(G(z)))$, on maximise $\log D(G(z))$ (<<~non-saturating loss~>>\index{non-saturating loss}).
\end{attention}

\section{Wasserstein GAN (WGAN)}
\label{sec:wgan}

\subsection{Distance de Earth Mover (Wasserstein-1)}

\begin{definition}[Distance de Wasserstein-1\index{distance de Wasserstein}\index{Earth Mover's distance}]
\begin{equation}
W(p_{\text{data}}, p_g) = \inf_{\gamma \in \Pi(p_{\text{data}}, p_g)} \E_{(x, y) \sim \gamma}\!\big[\|x - y\|\big]
\end{equation}
o\`u $\Pi(p_{\text{data}}, p_g)$ est l'ensemble des distributions jointes dont les marginales sont $p_{\text{data}}$ et $p_g$.
\end{definition}

\begin{theoreme}[Dualit\'e de Kantorovich-Rubinstein]
\begin{equation}
W(p_{\text{data}}, p_g) = \sup_{\|f\|_L \leq 1} \E_{x \sim p_{\text{data}}}\!\big[f(x)\big] - \E_{x \sim p_g}\!\big[f(x)\big]
\label{eq:kantorovich}
\end{equation}
o\`u le supremum est pris sur les fonctions 1-Lipschitz\index{Lipschitz}.
\end{theoreme}

\begin{intuition}
La distance de Wasserstein est plus douce que la JSD : m\^eme lorsque les supports de $p_{\text{data}}$ et $p_g$ ne se chevauchent pas, $W$ fournit un gradient exploitable, alors que $\KL$ et JSD sont ou infinis ou constants.
\end{intuition}

\begin{formules}[title=\textbf{Objectif WGAN}]\index{WGAN}
Le discriminateur devient un <<~critique~>> $f_w$ (sans sigmoid en sortie) :
\begin{equation}
\max_{w : \|f_w\|_L \leq 1} \; \E_{x \sim p_{\text{data}}}\!\big[f_w(x)\big] - \E_{z \sim p_z}\!\big[f_w(G_\theta(z))\big]
\label{eq:wgan-objective}
\end{equation}
\end{formules}

\subsection{WGAN-GP : p\'enalit\'e de gradient}
\label{sec:wgan-gp}

Arjovsky et al.\ imposent la contrainte de Lipschitz par \emph{weight clipping}, ce qui est probl\'ematique. Gulrajani et al.\ (2017) proposent une \textbf{p\'enalit\'e de gradient}\index{gradient penalty}\index{WGAN-GP} :

\begin{formules}[title=\textbf{WGAN-GP}]
\begin{equation}
\mathcal{L}_{\text{WGAN-GP}} = \underbrace{\E_{z}\!\big[f_w(G(z))\big] - \E_{x}\!\big[f_w(x)\big]}_{\text{objectif WGAN}} + \lambda \,\underbrace{\E_{\hat{x}}\!\big[\big(\|\nabla_{\hat{x}} f_w(\hat{x})\|_2 - 1\big)^2\big]}_{\text{p\'enalit\'e de gradient}}
\label{eq:wgan-gp}
\end{equation}
o\`u $\hat{x} = \epsilon x + (1 - \epsilon) G(z)$ avec $\epsilon \sim U(0, 1)$ et typiquement $\lambda = 10$.
\end{formules}

\section{GAN conditionnel (cGAN)}
\label{sec:cgan}

\begin{definition}[Conditional GAN\index{conditional GAN}\index{cGAN}]
On conditionne \`a la fois $G$ et $D$ sur une information auxiliaire $y$ (label de classe, texte, etc.) :
\begin{equation}
\min_G \max_D \; \E_{x, y}\!\big[\log D(x, y)\big] + \E_{z, y}\!\big[\log\big(1 - D(G(z, y), y)\big)\big]
\end{equation}
\end{definition}

\section{DCGAN : directives architecturales}
\label{sec:dcgan}

Le \textbf{Deep Convolutional GAN}\index{DCGAN} (Radford et al., 2016) \'etablit des r\`egles empiriques pour stabiliser l'entra\^inement des GAN convolutifs :

\begin{bonnepratique}
\begin{enumerate}
    \item Remplacer le pooling par des convolutions \`a pas (\emph{strided convolutions}) dans $D$ et des convolutions transpos\'ees dans $G$.
    \item Utiliser la \emph{Batch Normalization}\index{Batch Normalization} dans $G$ et $D$ (sauf la derni\`ere couche de $G$ et la premi\`ere couche de $D$).
    \item Supprimer les couches enti\`erement connect\'ees.
    \item Utiliser ReLU dans $G$ (sauf la sortie : $\tanh$) et LeakyReLU dans $D$.
\end{enumerate}
\end{bonnepratique}

\section{Impl\'ementation PyTorch : DCGAN sur Fashion-MNIST}
\label{sec:dcgan-pytorch}

\begin{codeblock}[title=\textbf{DCGAN complet avec boucle d'entra\^inement}]
\begin{minted}{python}
import torch
import torch.nn as nn
from torchvision import datasets, transforms
from torch.utils.data import DataLoader

# Hyperparametres
LATENT_DIM = 100
IMG_CHANNELS = 1
FEATURES_G = 64
FEATURES_D = 64
LR = 2e-4
BETAS = (0.5, 0.999)
BATCH_SIZE = 128
NUM_EPOCHS = 25

class Generator(nn.Module):
    """Generateur DCGAN pour images 28x28."""
    def __init__(self, latent_dim, features_g, img_channels):
        super().__init__()
        self.net = nn.Sequential(
            # Entree: (batch, latent_dim, 1, 1)
            nn.ConvTranspose2d(latent_dim, features_g * 4, 4, 1, 0,
                               bias=False),
            nn.BatchNorm2d(features_g * 4),
            nn.ReLU(True),
            # -> (batch, features_g*4, 4, 4)
            nn.ConvTranspose2d(features_g * 4, features_g * 2, 3, 2, 1,
                               bias=False),
            nn.BatchNorm2d(features_g * 2),
            nn.ReLU(True),
            # -> (batch, features_g*2, 7, 7)
            nn.ConvTranspose2d(features_g * 2, features_g, 4, 2, 1,
                               bias=False),
            nn.BatchNorm2d(features_g),
            nn.ReLU(True),
            # -> (batch, features_g, 14, 14)
            nn.ConvTranspose2d(features_g, img_channels, 4, 2, 1,
                               bias=False),
            nn.Tanh(),
            # -> (batch, img_channels, 28, 28)
        )

    def forward(self, z):
        return self.net(z)


class Discriminator(nn.Module):
    """Discriminateur DCGAN pour images 28x28."""
    def __init__(self, img_channels, features_d):
        super().__init__()
        self.net = nn.Sequential(
            # Entree: (batch, img_channels, 28, 28)
            nn.Conv2d(img_channels, features_d, 4, 2, 1, bias=False),
            nn.LeakyReLU(0.2, inplace=True),
            # -> (batch, features_d, 14, 14)
            nn.Conv2d(features_d, features_d * 2, 4, 2, 1, bias=False),
            nn.BatchNorm2d(features_d * 2),
            nn.LeakyReLU(0.2, inplace=True),
            # -> (batch, features_d*2, 7, 7)
            nn.Conv2d(features_d * 2, features_d * 4, 3, 2, 1, bias=False),
            nn.BatchNorm2d(features_d * 4),
            nn.LeakyReLU(0.2, inplace=True),
            # -> (batch, features_d*4, 4, 4)
            nn.Conv2d(features_d * 4, 1, 4, 1, 0, bias=False),
            nn.Sigmoid(),
            # -> (batch, 1, 1, 1)
        )

    def forward(self, x):
        return self.net(x).view(-1, 1)


def weights_init(m):
    """Initialisation des poids selon les recommandations DCGAN."""
    classname = m.__class__.__name__
    if 'Conv' in classname:
        nn.init.normal_(m.weight.data, 0.0, 0.02)
    elif 'BatchNorm' in classname:
        nn.init.normal_(m.weight.data, 1.0, 0.02)
        nn.init.constant_(m.bias.data, 0)


# --- Preparation des donnees ---
transform = transforms.Compose([
    transforms.ToTensor(),
    transforms.Normalize([0.5], [0.5]),  # -> [-1, 1]
])
dataset = datasets.FashionMNIST(
    root='./data', train=True, download=True, transform=transform
)
dataloader = DataLoader(dataset, batch_size=BATCH_SIZE, shuffle=True,
                        num_workers=2, drop_last=True)

# --- Initialisation ---
device = torch.device('cuda' if torch.cuda.is_available() else 'cpu')
G = Generator(LATENT_DIM, FEATURES_G, IMG_CHANNELS).to(device)
D = Discriminator(IMG_CHANNELS, FEATURES_D).to(device)
G.apply(weights_init)
D.apply(weights_init)

opt_G = torch.optim.Adam(G.parameters(), lr=LR, betas=BETAS)
opt_D = torch.optim.Adam(D.parameters(), lr=LR, betas=BETAS)
criterion = nn.BCELoss()

fixed_noise = torch.randn(64, LATENT_DIM, 1, 1, device=device)

# --- Boucle d'entrainement ---
for epoch in range(NUM_EPOCHS):
    for i, (real_imgs, _) in enumerate(dataloader):
        real_imgs = real_imgs.to(device)
        batch_size = real_imgs.size(0)
        real_labels = torch.ones(batch_size, 1, device=device)
        fake_labels = torch.zeros(batch_size, 1, device=device)

        # (1) Mise a jour du discriminateur
        noise = torch.randn(batch_size, LATENT_DIM, 1, 1, device=device)
        fake_imgs = G(noise).detach()

        loss_D_real = criterion(D(real_imgs), real_labels)
        loss_D_fake = criterion(D(fake_imgs), fake_labels)
        loss_D = loss_D_real + loss_D_fake

        opt_D.zero_grad()
        loss_D.backward()
        opt_D.step()

        # (2) Mise a jour du generateur
        noise = torch.randn(batch_size, LATENT_DIM, 1, 1, device=device)
        fake_imgs = G(noise)
        loss_G = criterion(D(fake_imgs), real_labels)  # non-saturating

        opt_G.zero_grad()
        loss_G.backward()
        opt_G.step()

    print(f"Epoch [{epoch+1}/{NUM_EPOCHS}]  "
          f"Loss_D: {loss_D.item():.4f}  Loss_G: {loss_G.item():.4f}")
\end{minted}
\end{codeblock}

\begin{codeoutput}
\begin{verbatim}
Epoch [1/25]   Loss_D: 0.5823  Loss_G: 1.8932
Epoch [2/25]   Loss_D: 0.4517  Loss_G: 2.1045
...
Epoch [10/25]  Loss_D: 0.6102  Loss_G: 1.3287
...
Epoch [25/25]  Loss_D: 0.6731  Loss_G: 0.9845
\end{verbatim}
\end{codeoutput}

\section{M\'etriques d'\'evaluation}
\label{sec:gan-metrics}

\subsection{Inception Score (IS)}

\begin{formules}[title=\textbf{Inception Score}]\index{Inception Score}\index{IS}
\begin{equation}
\text{IS} = \exp\!\Big(\E_{x \sim p_g}\!\big[\KL\!\big(p(y \mid x) \,\|\, p(y)\big)\big]\Big)
\end{equation}
o\`u $p(y \mid x)$ est donn\'e par un r\'eseau Inception pr\'e-entra\^in\'e et $p(y) = \E_x[p(y \mid x)]$.

Un IS \'elev\'e signifie : (1) les images sont nettes ($p(y|x)$ concentr\'ee) et (2) les images sont vari\'ees ($p(y)$ uniforme).
\end{formules}

\subsection{Fr\'echet Inception Distance (FID)}

\begin{formules}[title=\textbf{Fr\'echet Inception Distance}]\index{FID}\index{Fr\'echet Inception Distance}
On extrait les features de la couche pool3 du r\'eseau Inception pour les vraies donn\'ees ($\mu_r, \Sigma_r$) et les donn\'ees g\'en\'er\'ees ($\mu_g, \Sigma_g$) :
\begin{equation}
\text{FID} = \|\mu_r - \mu_g\|_2^2 + \text{Tr}\!\Big(\Sigma_r + \Sigma_g - 2\big(\Sigma_r \Sigma_g\big)^{1/2}\Big)
\label{eq:fid}
\end{equation}
Un FID \textbf{bas} indique une meilleure qualit\'e et diversit\'e.
\end{formules}

\begin{remarque}
La FID est pr\'ef\'er\'ee \`a l'IS car elle compare la distribution g\'en\'er\'ee aux \emph{vraies} donn\'ees, tandis que l'IS n'\'evalue que les propri\'et\'es internes des images g\'en\'er\'ees.
\end{remarque}

\section{\'Etude d'ablation}
\label{sec:gan-ablation}

\begin{table}[htbp]
\centering
\caption{\'Etude d'ablation DCGAN sur Fashion-MNIST. FID mesur\'e apr\`es 25 \'epoques (plus bas = meilleur).}
\label{tab:gan-ablation}
\begin{tabular}{lccc}
\toprule
\textbf{Configuration} & \textbf{Dim.\ latente} & \textbf{LR ($G$ / $D$)} & \textbf{FID} $\downarrow$ \\
\midrule
Base & 100 & $2\!\times\!10^{-4}$ / $2\!\times\!10^{-4}$ & 28.3 \\
\midrule
$z \in \R^{32}$ & 32 & $2\!\times\!10^{-4}$ / $2\!\times\!10^{-4}$ & 35.1 \\
$z \in \R^{64}$ & 64 & $2\!\times\!10^{-4}$ / $2\!\times\!10^{-4}$ & 30.2 \\
$z \in \R^{256}$ & 256 & $2\!\times\!10^{-4}$ / $2\!\times\!10^{-4}$ & 27.8 \\
\midrule
LR $G$ plus \'elev\'e & 100 & $5\!\times\!10^{-4}$ / $2\!\times\!10^{-4}$ & 32.5 \\
LR $D$ plus \'elev\'e & 100 & $2\!\times\!10^{-4}$ / $5\!\times\!10^{-4}$ & 41.7 \\
TTUR ($G < D$) & 100 & $1\!\times\!10^{-4}$ / $4\!\times\!10^{-4}$ & 25.9 \\
\midrule
$n_{\text{critic}}=1$ (standard) & 100 & $2\!\times\!10^{-4}$ / $2\!\times\!10^{-4}$ & 28.3 \\
$n_{\text{critic}}=3$ & 100 & $2\!\times\!10^{-4}$ / $2\!\times\!10^{-4}$ & 26.7 \\
$n_{\text{critic}}=5$ (WGAN-GP) & 100 & $1\!\times\!10^{-4}$ / $1\!\times\!10^{-4}$ & 23.4 \\
\bottomrule
\end{tabular}
\end{table}

\begin{remarque}
Points cl\'es : (1) la dimension latente a un impact mod\'er\'e, mais trop petite elle limite la diversit\'e ; (2) un taux d'apprentissage plus \'elev\'e pour $D$ (TTUR --- Two Time-scale Update Rule)\index{TTUR} am\'eliore la FID ; (3) WGAN-GP avec $n_{\text{critic}}=5$ offre la meilleure performance.
\end{remarque}

\section{\'Etat de l'art et connexions}
\label{sec:gan-sota}

\begin{sota}[title=\textbf{Au-del\`a des GAN classiques}]
\begin{itemize}
    \item \textbf{StyleGAN}\index{StyleGAN} (Karras et al., 2019--2021) : introduit un r\'eseau de mapping $z \to w$ et l'injection de style via \emph{Adaptive Instance Normalization} (AdaIN). StyleGAN3 \'elimine les artefacts de texture par une \'equivariance stricte aux translations.

    \item \textbf{Mod\`eles de diffusion}\index{mod\`eles de diffusion} (Ho et al., 2020 ; Dhariwal et Nichol, 2021) : les mod\`eles de d\'ebruitage diffusif ont largement surpass\'e les GAN en termes de FID sur ImageNet, tout en offrant une diversit\'e sup\'erieure et un entra\^inement plus stable. Voir DALL-E 2, Stable Diffusion, Imagen.

    \item \textbf{Consistency Models}\index{Consistency Models} (Song et al., 2023) : distillent les mod\`eles de diffusion pour permettre la g\'en\'eration en une seule \'etape, combinant la qualit\'e de la diffusion avec la rapidit\'e des GAN.

    \item Les GAN restent pr\'ef\'er\'es dans des applications n\'ecessitant une g\'en\'eration ultra-rapide (super-r\'esolution en temps r\'eel, synth\`ese vid\'eo).
\end{itemize}
\end{sota}

\section{Exercices}
\label{sec:gan-exercices}

\begin{exercice}[V\'erification du discriminateur optimal]
\label{exo:optimal-disc}
\textbf{Difficult\'e~: 1/3}
Soit $p_{\text{data}} = \mathcal{N}(3, 1)$ et $p_g = \mathcal{N}(0, 1)$. Calculez analytiquement $D^*(x)$ en utilisant l'\'equation~\eqref{eq:optimal-D}. Tracez $D^*(x)$ pour $x \in [-5, 8]$. Commentez la forme de la courbe.
\end{exercice}

\begin{exercice}[WGAN : comparaison avec GAN standard]
\label{exo:wgan-compare}
\textbf{Difficult\'e~: 2/3}
Consid\'erons deux distributions : $p_{\text{data}} = \delta_0$ (masse de Dirac en 0) et $p_g = \delta_\theta$.
\begin{enumerate}
    \item Calculez $\text{JSD}(p_{\text{data}} \| p_g)$ pour $\theta \neq 0$ et $\theta = 0$.
    \item Calculez $W(p_{\text{data}}, p_g) = |\theta|$.
    \item En d\'eduire pourquoi $W$ fournit un gradient utile pour $G$ tandis que JSD ne le fait pas.
\end{enumerate}
\end{exercice}

\begin{exercice}[Impl\'ementation WGAN-GP]
\label{exo:wgan-gp-impl}
\textbf{Difficult\'e~: 2/3}
Modifiez le code DCGAN de la section~\ref{sec:dcgan-pytorch} pour :
\begin{enumerate}
    \item Supprimer le sigmoid du discriminateur (<<~critique~>>).
    \item Impl\'ementer la p\'enalit\'e de gradient (\'equation~\eqref{eq:wgan-gp}).
    \item Entra\^iner le critique $n_{\text{critic}}=5$ fois par it\'eration du g\'en\'erateur.
    \item Comparer les courbes de perte et la qualit\'e visuelle avec le GAN standard.
\end{enumerate}
\end{exercice}

\begin{exercice}[Conditional GAN pour la g\'en\'eration cibl\'ee]
\label{exo:cgan-impl}
\textbf{Difficult\'e~: 3/3}
Impl\'ementez un cGAN\index{cGAN} sur Fashion-MNIST qui prend en entr\'ee un label de classe et g\'en\`ere une image de la cat\'egorie correspondante.

\emph{Indice} : concat\'enez un \emph{one-hot encoding} du label au vecteur latent pour $G$, et ajoutez un \emph{embedding} du label comme canal suppl\'ementaire pour $D$.
\end{exercice}

\begin{exercice}[Analyse th\'eorique --- convergence des GAN]
\label{exo:gan-convergence}
\textbf{Difficult\'e~: 3/3}
\begin{enumerate}
    \item Montrez que si $D$ et $G$ sont mis \`a jour simultan\'ement par descente de gradient sur $V(D, G)$, le syst\`eme dynamique n'est \emph{pas} un jeu \`a gradient (les gradients ne sont pas conservatifs). Que cela implique-t-il pour la convergence~?
    \item Montrez que pour le GAN standard en dimension 1, avec $p_{\text{data}} = \mathcal{N}(0, 1)$ et $G(z) = \mu + \sigma z$ ($z \sim \mathcal{N}(0,1)$), le point fixe $\mu=0, \sigma=1$ n'est pas un attracteur stable du syst\`eme dynamique de gradient simultan\'e.
\end{enumerate}
\end{exercice}

\begin{exercice}[Calcul de la FID {\normalfont (projet)}]
\label{exo:fid-compute}
\textbf{Difficult\'e~: 3/3}
\begin{enumerate}
    \item Impl\'ementez le calcul de la FID en utilisant un r\'eseau pr\'e-entra\^in\'e (e.g., InceptionV3 ou un r\'eseau plus simple pour des images $28 \times 28$).
    \item Calculez la FID de votre DCGAN \`a diff\'erentes \'epoques d'entra\^inement.
    \item Tracez l'\'evolution de la FID en fonction de l'\'epoque et commentez.
\end{enumerate}
\end{exercice}
